% §11.6 Bibliographic Notes
%
% Cross-reference resolution status:
%
% ALREADY RESOLVED (within chapter 11, drafted in Tasks 2-6 of Plan 11):
%   \Cref{sec:miller-rabin-formal}                -- §11.1
%   \Cref{thm:miller-rabin-formal}                -- §11.1 theorem
%   \Cref{sec:las-vegas-vs-monte-carlo}           -- §11.2
%   \Cref{def:las-vegas}                          -- §11.2
%   \Cref{def:monte-carlo}                        -- §11.2
%   \Cref{sec:rp-bpp}                             -- §11.3
%   \Cref{def:rp}                                 -- §11.3
%   \Cref{def:bpp}                                -- §11.3
%   \Cref{sec:sample-complexity}                  -- §11.4
%   \Cref{thm:chernoff}                           -- §11.4 theorem
%   \Cref{sec:ch11-bridge}                        -- §11.5
%
% ALREADY RESOLVED (earlier chapters):
%   \Cref{sec:miller-rabin}                       -- chapter 3 §3.4 (informal walkthrough)
%   \Cref{sec:binomial-estimator}                 -- chapter 7 §7.2
%
% Citations resolved by existing references.bib entries:
%   \cite{miller-1976}                            -- Miller (1976), already present
%   \cite{rabin-1980}                             -- Rabin (1980), already present
%   \cite{mitzenmacher-upfal}                     -- Mitzenmacher-Upfal, already present
%
% Citations resolved by new references.bib entries added in this task:
%   \cite{motwani-raghavan}                       -- Motwani-Raghavan (1995)
%   \cite{karger-1993}                            -- Karger (1993) min-cut

\section*{Bibliographic Notes}
\addcontentsline{toc}{section}{Bibliographic Notes}
\label{sec:bib-notes-ch11}

\paragraph{Miller-Rabin primality testing.}
The formal $\Bernoulli{\BoolType}$ statement of
\Cref{thm:miller-rabin-formal} packages two ingredients of separate
origin. The per-round bound $\fprate \le 1/4$ is
\textcite{rabin-1980}'s strong-liar theorem, which sharpens to
$1/4$ a polynomial-time randomized primality test that
\textcite{miller-1976} had earlier shown to be deterministic and
correct under the generalized Riemann hypothesis. Rabin's
contribution was to replace the GRH assumption with random sampling
of bases, trading determinism for an unconditional probabilistic
guarantee and a $1/4$ rather than $1/2$ per-round false-positive
rate. The number-theoretic content of \Cref{sec:miller-rabin-formal}
is theirs; the framework reading, in which the witness predicate is
a $\Bernoulli{\BoolType}$ and the $k$-round bound
$\fprate^{(k)} \le (1/4)^k$ is a corollary of the composition
theorem of chapter~3 applied to independent samples under chapter~4
Axiom~1, is the chapter's specialization rather than theirs. A
historical and pedagogical comparison of the deterministic
Miller test and the probabilistic Rabin test is given in chapter~31
of \textcite{cormen-clrs}; the AKS deterministic polynomial-time
primality test of Agrawal, Kayal, and Saxena (2002) supersedes both
in worst-case complexity terms but is not competitive in
constant factors with $k$-round Miller-Rabin for the input sizes
that arise in practical cryptographic key generation, which is why
Miller-Rabin remains the standard library implementation.

\paragraph{Randomized algorithms as a field.}
The standard graduate-level reference on randomized algorithms is
\textcite{motwani-raghavan}, whose 1995 Cambridge textbook covers
the Las~Vegas/Monte~Carlo distinction of
\Cref{sec:las-vegas-vs-monte-carlo}, the RP/BPP/ZPP classification of
\Cref{sec:rp-bpp}, the Chernoff and Hoeffding tail bounds of
\Cref{sec:sample-complexity}, and dozens of specific algorithms
(randomized quicksort, Karger min-cut, polynomial identity testing,
pattern matching, geometric algorithms, approximate counting, Markov
chain Monte Carlo) that the chapter treats only by example or not at
all. Motwani-Raghavan organizes its material by technique;
the chapter organizes by error structure. The two organizations are
complementary: a reader who wants to apply randomization to a new
problem should treat Motwani-Raghavan as the next book, and a reader
who has finished Motwani-Raghavan and wants the unifying type-algebra
view should treat this book as the complement.

\paragraph{Chernoff bounds and concentration.}
The multiplicative Chernoff bound of \Cref{thm:chernoff} and the
moment-generating-function technique that gives it originated in
Chernoff's 1952 \emph{Annals of Mathematical Statistics} paper on
the efficiency of tests based on sums of observations; the textbook
formulation used in \Cref{sec:sample-complexity} and the BPP
amplification calculation follow chapter~4 of
\textcite{mitzenmacher-upfal}, the recommended companion to the
chapter for any reader who wants the concentration-inequality
machinery developed more fully. Mitzenmacher and Upfal also give the
additive Hoeffding inequality, the Azuma-Hoeffding martingale
generalization, and McDiarmid's bounded-difference inequality, all
of which subsume \Cref{thm:chernoff} in different directions and are
needed for randomized algorithms whose error is not a simple sum of
independent Bernoulli trials. The chapter uses only the
multiplicative form because that is what the Miller-Rabin and BPP
applications need; in chapter~10 the additive form underwrites the
CountMin and HyperLogLog error bounds. The choice between additive
and multiplicative formulations is problem-driven rather than a
matter of mathematical taste.

\paragraph{Karger's min-cut algorithm.}
The randomized min-cut algorithm cited in
\Cref{sec:las-vegas-vs-monte-carlo} as the canonical Monte Carlo
example beyond Miller-Rabin is due to \textcite{karger-1993}, whose
SODA paper introduces the contraction algorithm: repeatedly pick a
uniformly random edge and contract its endpoints, terminate when two
vertices remain, and return the cut between them. The analysis
shows that any fixed min-cut survives all contractions with
probability at least $2/(n(n-1))$, so $\Theta(n^2 \log n)$
independent trials suffice to find a min-cut with high probability,
in $\widetilde O(n^4)$ total time. Karger's later refinements with
Stein and others bring this to $\widetilde O(n^2)$ via the recursive
contraction tree, which is the asymptotically fastest known
algorithm for the unweighted min-cut problem. The pedagogical value
of Karger's algorithm in the present chapter is not its complexity
result but its illustration of the Monte Carlo design pattern: a
fixed-runtime randomized procedure with a small per-trial success
probability is converted to a high-confidence algorithm by
amplification through independent repetition, which is the same
$\Bernoulli{\BoolType}$ amplification analyzed in
\Cref{sec:sample-complexity} for Miller-Rabin and BPP, applied to a
combinatorial optimization problem rather than a decision problem.
